% Target: rmp-iii-b-braid-two
% Formula: The second braid panel exchanges the a and b strings around i and j.
% Ink: Canonical public kernel interlaced arcs and four marked anyon sites.
% Boundary: Four lattice legs and six string endpoints are open.
% Source: paper TN-Review-main.tex line 1676; author source ImagesReview Section
%   III/ImagesReview Section III.tex lines 826-862
% Capabilities: kernel, multi-strand-braid, crossing-order
\[
  % The author's panel weaves the two strings through the plaquette without a
  % string-string crossing: the b string rounds the anyon i across its south
  % leg, and each string-lattice crossing carries its dot.  The kernel routes
  % below say exactly that; every dot stands at the address of its crossing.
  \tenkzkernel
  \tndeclare{species}{operator}{hue=source:red}
  \begin{tenkz}[rows={wire,wire}, cols=2, bonds=none]
    \tn[at=(1,1), name=b, role=extra, label pos=nw]{b}
    \tn[at=(1,2), name=a, role=marked, label pos=ne]{a}
    \tn[at=(2,1), name=j, role=extra, label pos=sw]{j}
    \tn[at=(2,2), name=i, role=marked, label pos=se]{i}
    \tnwire{open w}{b.w} \tnwire{b.e}{a.w} \tnwire{a.e}{open e}
    \tnwire{open w}{j.w} \tnwire{j.e}{i.w} \tnwire{i.e}{open e}
    \tnwire{b.n}{open n} \tnwire{b.s}{j.n} \tnwire{j.s}{open s}
    \tnwire{a.n}{open n} \tnwire{a.s}{i.n} \tnwire{i.s}{open s}
    % Both strings are junction chains (retired via=): a named marker bead at
    % each waypoint that used to be dotted (the same records the source
    % marks its crossings with), plus one invisible junction at sb's sixth
    % waypoint, which the original left undotted because it is a routing
    % point, not a crossing.  Every consecutive pair of segments also needs
    % its own declared order at the corner they share, same as the
    % torus-two wrap chain (#6897).
    \tn[at=0.3 w of b, name=m1]{}
    \tn[at=midway b and j, name=m2]{}
    \tn[at=midway j and i, name=m3]{}
    \tn[at=0.3 s of i, name=m4]{}
    \tn[at=0.3 e of i, name=m5]{}
    \tn[at=0.3 e of midway i and a, skin=none, name=m6]{}
    \tn[at=midway i and a, name=n1]{}
    \tn[at=midway b and a, name=n2]{}
    \tn[at=0.3 n of a, name=n3]{}
    \tnwire[kind=string, species=operator, name=sb1]{open n}{m1}
    \tnwire[kind=string, species=operator, name=sb2,
            cross={over at crossing of sb1 and sb2}]{m1}{m2}
    \tnwire[kind=string, species=operator, name=sb3,
            cross={over at crossing of sb2 and sb3}]{m2}{m3}
    \tnwire[kind=string, species=operator, name=sb4,
            cross={over at crossing of sb3 and sb4}]{m3}{m4}
    \tnwire[kind=string, species=operator, name=sb5,
            cross={over at crossing of sb4 and sb5}]{m4}{m5}
    \tnwire[kind=string, species=operator, name=sb6,
            cross={over at crossing of sb5 and sb6}]{m5}{m6}
    \tnwire[kind=string, species=operator, name=sb7,
            cross={over at crossing of sb6 and sb7}]{m6}{a.315}
    \tnwire[kind=string, species=operator, name=sa1]{i.45}{n1}
    \tnwire[kind=string, species=operator, name=sa2,
            cross={over at crossing of sa1 and sa2}]{n1}{n2}
    \tnwire[kind=string, species=operator, name=sa3,
            cross={over at crossing of sa2 and sa3}]{n2}{n3}
    \tnwire[kind=string, species=operator, name=sa4,
            cross={over at crossing of sa3 and sa4}]{n3}{open e}
  \end{tenkz}
\]
